% !TEX root = ../main.tex
\section{Variants}
Since its introduction, ROCKET has inspired a family of successors that address its main limitations—computational overhead, non-determinism, and restriction to univariate input—while preserving or improving upon its accuracy. The three principal variants are described below.

\subsection{MiniROCKET}
MiniROCKET (Dempster et al., 2021), optimizes the original framework by over an order of magnitude in computational speed without compromising classification accuracy. This variant transitions the model toward a nearly deterministic state by eliminating several stochastic components.
In MiniROCKET, the kernel length is standardized to a fixed value of 9, and the weight space is constrained to only two discrete values. Unlike the original random initialization, the bias values are systematically derived from the distribution of the convolution outputs on the training data. Furthermore, the feature set is streamlined by discarding the "max" pooling operation, extracting only the Proportion of Positive Values (PPV). These structural constraints enable hardware-level optimizations that exploit the fixed parameters, providing a substantial reduction in training and inference latency while maintaining state-of-the-art performance across the UCR archive.

\subsection{MultiROCKET}
MultiROCKET (Tan et al., 2022), extends the framework to multivariate time series and introduces a significantly richer feature vocabulary. While maintaining the efficient weight scheme of MiniROCKET (utilizing the fixed pool of 84 binary vectors), MultiROCKET augments the feature extraction process by introducing three additional pooling operators beyond the original PPV and Max:
\begin{itemize}
	\item \textbf{Mean of Positive Values (MPV)}: Captures the average activation strength, quantifying the intensity of a pattern match.
	\item \textbf{Mean of Indices of Positive Values (MIPV)}: Encodes temporal localization by identifying where in the series a pattern typically occurs.
	\item \textbf{Longest Stretch of Positive Values (LSPV)}: Measures the temporal persistence of a pattern through its longest contiguous activation.
\end{itemize}
These five operators produce $5\times k$ features per channel, providing the linear classifier with critical information regarding the frequency, strength, location, and duration of detected patterns. For multivariate data with $C$ channels, the transformation is applied independently to each channel and concatenated. This approach achieves a statistically significant improvement in mean rank across the UCR archive while preserving a linear computational complexity of $O(k\cdot n\cdot l\cdot C)$. By avoiding cross-channel combinatorial kernels, MultiROCKET remains highly scalable while capturing a more comprehensive representation of complex temporal dynamics.

\subsection{Hydra}
Hydra (Dempster et al., 2023) represents a conceptually distinct evolution of the ROCKET paradigm, combining ideas from random convolutional kernels with dictionary-based methods.

Rather than summarising each feature map with scalar pooling statistics, Hydra organises kernels into groups of $g$ kernels sharing the same dilation and padding, and builds a histogram of nearest-kernel assignments across each group and each position in the series. This histogram --- counting, for each kernel within a group, how many time-series positions are most similar to that kernel's weights --- serves as the feature representation passed to the linear classifier.

This construction is closely related to the Bag-of-SFA-Symbols (BOSS) dictionary representation: the group of $g$ kernels can be seen as a soft codebook, and the histogram as a symbolic frequency profile. The key difference is that Hydra's codebook is random and extremely cheap to compute, whereas BOSS requires a Symbolic Fourier Approximation discretisation step that is quadratic in training-set size. Empirically, Hydra surpasses ROCKET and MultiROCKET in mean classification rank on the UCR archive while remaining an order of magnitude faster than InceptionTime. It is particularly effective on problems with long and complex series, where the positional frequency of matched patterns --- rather than just their peak value or overall proportion --- carries the most discriminative information.

A combined variant, \textbf{Hydra-MultiROCKET}, concatenates Hydra's histogram features with MultiROCKET's pooling features, achieving the highest mean accuracy in the ROCKET family to date on the UCR archive.
